% ==============================================================================
% ABSTRACT & KURZFASSUNG
% Chair of Wireless Systems (Fachgebiet Drahtlose Systeme)
% Brandenburg University of Technology Cottbus-Senftenberg (BTU)
% ==============================================================================

% ------------------------------------------------------------------------------
% Student Writing Guide: Abstract
% ------------------------------------------------------------------------------
\begin{tcolorbox}[
    colback=white,
    colframe=btuRed,
    title={\textbf{Student Writing Guide: Crafting an Executive Abstract}},
    fonttitle=\bfseries\sffamily,
    sharp corners,
    boxrule=0.8pt
]
\small
The abstract is the most frequently read part of your thesis. It must stand independently without referring to the main text, figures, or literature citations. Follow the proven \textbf{5-sentence structural model} (target: 150--250 words):
\begin{enumerate}[leftmargin=*,itemsep=1pt,topsep=2pt]
    \item \textbf{Context \& Motivation (1--2 sentences):} What domain are you addressing (e.g., edge intelligence, energy limits in WSNs) and why is it important?
    \item \textbf{Problem Statement \& Gap (1--2 sentences):} What specific challenge remains unsolved by existing systems (e.g., monolithic uncertainty, rigid task offloading)?
    \item \textbf{Proposed Methodology (2--3 sentences):} What concrete architecture, algorithm, or protocol do you propose (e.g., adaptive agentic orchestration, uncertainty-aware offloading)?
    \item \textbf{Empirical Results (2--3 sentences):} What quantitative results did your testbed or benchmark achieve (e.g., reduced latency by 27\%, saved energy by 34\%)?
    \item \textbf{Scientific Contribution \& Impact (1 sentence):} What is the primary takeaway for the research community and industry?
\end{enumerate}
\textit{Note: Both an English Abstract and a German Kurzfassung are required at BTU Cottbus--Senftenberg.}
\end{tcolorbox}

\vspace{0.4cm}

% ------------------------------------------------------------------------------
% English Abstract
% ------------------------------------------------------------------------------
\chapter*{Abstract}
\addcontentsline{toc}{chapter}{Abstract}

Modern distributed cyber-physical environments, wireless sensor networks (WSNs), and low-power edge computing architectures increasingly rely on autonomous decision-making agents to process high-dimensional telemetry locally while offloading computationally intensive inference to remote servers. However, existing orchestration frameworks treat uncertainty as a monolithic scalar quantity, confounding \emph{epistemic uncertainty} (lack of model knowledge, resolvable by retrieval or cloud escalation) with \emph{aleatoric uncertainty} (inherent channel noise, sensor drift, and stochastic environmental dynamics). Consequently, conventional task dispatchers execute sub-optimal offloading decisions: depleting constrained battery reserves on fruitless cloud escalations or failing to intervene when critical sensor anomalies occur.

This thesis presents an adaptive framework for \emph{Adaptive Agentic Orchestration and Uncertainty-Aware Task Offloading} tailored to resource-constrained wireless sensor nodes and edge gateways. By decomposing overall decision uncertainty into orthogonal epistemic and aleatoric components via semantic entropy clustering and lightweight natural language inference, the proposed architecture dynamically decides whether to execute subtasks locally, offload them across energy-efficient wireless links, or initiate structured human-in-the-loop escalation. We evaluate the proposed architecture using standard reasoning benchmarks as well as a physical wireless sensor testbed comprising battery-operated edge nodes. Empirical results demonstrate that the proposed system reduces unnecessary wireless transmissions by up to \SI{38}{\percent} and lowers end-to-end inference latency by \SI{27}{\percent} compared to confidence-agnostic offloading baselines, while extending wireless node operational lifetime by \SI{31}{\percent}. The findings establish that fine-grained uncertainty decomposition is a foundational requirement for dependable, energy-efficient, and resilient edge intelligence in next-generation wireless systems.

\vspace{0.8cm}
\noindent\textbf{Keywords:} \ThesisKeywords.

\vfill
\newpage

% ------------------------------------------------------------------------------
% Deutsche Kurzfassung
% ------------------------------------------------------------------------------
\chapter*{Kurzfassung}
\addcontentsline{toc}{chapter}{Kurzfassung}

Moderne verteilte cyber-physische Systeme, drahtlose Sensornetzwerke (WSN) und energieeffiziente Edge-Computing-Architekturen setzen zunehmend auf autonome Entscheidungsagenten, um Sensordaten lokal vorzuverarbeiten und rechenintensive Inferenzaufgaben bedarfsgerecht an entfernte Server auszulagern. Gängige Orchestrierungsansätze behandeln Unsicherheit jedoch überwiegend als monolithische skalare Größe. Hierbei wird \emph{epistemische Unsicherheit} (Wissensmangel des Modells, behebbar durch externe Wissensabfragen oder Cloud-Eskalation) mit \emph{aleatorischer Unsicherheit} (inhärentes Kanalrauschen, Sensordrift oder stochastische Umweltbedingungen) vermischt. Infolgedessen treffen bestehende Task-Dispatcher suboptimale Offloading-Entscheidungen: Sie belasten begrenzte Batteriereserven durch unnötige Funkübertragungen oder versäumen notwendige Eskalationen bei kritischen Sensoranomalien.

In dieser Masterarbeit wird eine Architektur für \emph{adaptive agentenbasierte Orchestrierung und unsicherheitsbewusste Aufgabenübertragung} (Task Offloading) für ressourcenbeschränkte drahtlose Sensorknoten vorgestellt. Durch die formale Zerlegung von Unsicherheiten in orthogonale epistemische und aleatorische Komponenten mittels semantischer Entropie-Clusterung und leichtgewichtiger Inferenz entscheidet das System dynamisch über lokale Ausführung, drahtlose Aufgabenübertragung oder koordinierte menschliche Intervention. Die experimentelle Validierung anhand etablierter Benchmarks sowie eines realen drahtlosen Sensortestbeds belegt eine Reduktion unnötiger Funkübertragungen um bis zu \SI{38}{\percent} und eine Verkürzung der P95-Latenz um \SI{27}{\percent} gegenüber referenzbasierten Verfahren bei einer Steigerung der Knotenlebensdauer um \SI{31}{\percent}. Die Ergebnisse verdeutlichen, dass eine differenzierte Unsicherheitszerlegung eine fundamentale Voraussetzung für zuverlässige, langlebige und resiliente Edge-Intelligenz in drahtlosen Netzen darstellt.

\vspace{0.8cm}
\begin{flushleft}
\noindent\textbf{Schlagwörter:} Drahtlose Sensornetzwerke, Edge-Intelligenz, Agentenbasierte Orchestrierung, Task Offloading, Unsicherheitsquantifizierung, Gemeinsame Kognitive Systeme.
\end{flushleft}
